\chapter{Cài đặt và triển khai hệ thống}

% ============================================================
% 5.1 Cài đặt môi trường phát triển
% ============================================================
\section{Cài đặt môi trường phát triển}

\subsection{Môi trường phát triển phần cứng}
% Để triển khai hệ thống IoT giám sát vận chuyển đạn dược, nhóm sử dụng nhiều bộ kit
% phát triển phần cứng khác nhau, đảm nhiệm các vai trò riêng trong kiến trúc tổng thể.
% Cụ thể, hai bộ kit LoRa Ra-08H tích hợp SoC ASR6601 được sử dụng lần lượt làm Node
% truyền (TX) và Gateway nhận (RX), đảm nhiệm chức năng truyền thông không dây
% tầm xa theo giao thức LoRa P2P. Vi điều khiển trung tâm của hệ thống là kit phát triển
% ESP32 NodeMCU, có nhiệm vụ thu thập dữ liệu từ các cảm biến, xử lý sơ bộ và
% truyền dữ liệu tới Gateway thông qua giao tiếp UART kết hợp với kết nối WiFi.

% Trong quá trình nạp chương trình và kiểm thử hệ thống, nhóm sử dụng mạch chuyển đổi
% USB-to-com để nạp firmware và theo dõi log debug của module
% Ra-08H, đồng thời sử dụng cáp Micro-USB để nạp chương trình và cấp nguồn cho
% vi điều khiển ESP32. Việc lựa chọn các công cụ này giúp quá trình phát triển,
% gỡ lỗi và triển khai hệ thống diễn ra thuận tiện và ổn định.
\subsubsection{Môi trường phát triển firmware LoRa P2P}
Do module Ra-08H sử dụng SoC ASR6601 dựa trên kiến trúc ARM Cortex-M4, quá trình phát
triển firmware cho hệ thống LoRa P2P yêu cầu bộ công cụ chuyên biệt do nhà sản xuất
cung cấp. Trong quá trình hiện thực, nhóm sử dụng môi trường phát triển Keil
$\mu$Vision~5 (MDK-ARM) để biên dịch và quản lý mã nguồn ngôn ngữ C cho vi điều khiển.
Đây là công cụ phổ biến trong phát triển hệ thống nhúng, hỗ trợ tốt cho kiến trúc
ARM và tích hợp đầy đủ các chức năng build, debug và quản lý bộ nhớ.

Bên cạnh đó, ASR6601 SDK được sử dụng làm nền tảng phát triển firmware, cung cấp
các thư viện driver phần cứng, cấu hình radio LoRa và các ví dụ mẫu, trong đó có
chương trình \texttt{pingpong} được nhóm lựa chọn làm cơ sở để phát triển giao thức
truyền thông LoRa P2P. Việc sử dụng SDK chính hãng giúp đảm bảo tính tương thích,
ổn định và rút ngắn thời gian triển khai hệ thống.

Để thuận tiện cho việc tạo và cấu hình dự án, nhóm đã sử dụng công cụ hỗ trợ mang tên
\texttt{KeilProjectGen}, là một script hỗ trợ sinh tự động file dự án
\texttt{.uvprojx} từ cấu trúc mã nguồn của SDK. Công cụ này giúp giảm thiểu các lỗi
cấu hình thủ công và đảm bảo sự đồng nhất giữa các project firmware.
Cuối cùng, quá trình nạp firmware vào bộ nhớ flash của chip ASR6601 được thực hiện
thông qua phần mềm TremoProg phiên bản~1.0, cho phép nạp file binary trực tiếp
vào module Ra-08H một cách nhanh chóng và ổn định.


\subsection{Môi trường phát triển phần mềm}
Môi trường phát triển phần mềm của hệ thống được xây dựng nhằm đáp ứng đồng thời các yêu cầu về lập trình vi điều khiển, triển khai backend xử lý dữ liệu và phát triển ứng dụng di động. 

\begin{itemize}
    \item \textbf{Vi điều khiển (ESP32):} Nhóm sử dụng Visual Studio Code kết hợp với Extension PlatformIO trên nền tảng Framework Arduino. Cách tiếp cận này giúp đơn giản hóa quá trình quản lý thư viện, biên dịch và nạp chương trình, đồng thời tăng khả năng mở rộng và bảo trì mã nguồn.
    \item \textbf{Backend \& Middleware:} Ngôn ngữ lập trình Python phiên bản 3.x được sử dụng để xây dựng script trung gian (Bridge) thực hiện chức năng kết nối giữa MQTT Broker và cơ sở dữ liệu Firebase. Việc lựa chọn Python cho phép phát triển nhanh, dễ dàng xử lý dữ liệu dạng JSON và tích hợp hiệu quả với các thư viện hỗ trợ như \texttt{firebase-admin}. Công tác quản lý và giám sát dữ liệu được thực hiện thông qua Google Firebase Console, đóng vai trò là nền tảng lưu trữ và đồng bộ dữ liệu thời gian thực cho hệ thống.
    
    \item \textbf{Mobile Application:} Nhóm sử dụng Flutter SDK phiên bản mới nhất để phát triển ứng dụng đa nền tảng, kết hợp với Android Studio nhằm hỗ trợ giả lập, kiểm thử và đóng gói ứng dụng dưới dạng file APK. Ngoài ra, hệ thống sử dụng phần mềm Mosquitto MQTT được cài đặt trên hệ điều hành Windows để làm MQTT Broker, chịu trách nhiệm quản lý kết nối, phân phối và truyền tải dữ liệu giữa các thành phần trong hệ thống.
\end{itemize}

% ============================================================
% 5.2 Cài đặt firmware cho Node Lora P2P
% ============================================================
\section{Cài đặt firmware cho Node Lora P2P}

\subsection{Các bước build firmware LoRa P2P}
Quy trình biên dịch và nạp firmware cho module Ra-08H được thực hiện theo một chuỗi
các bước tuần tự nhằm đảm bảo firmware được cấu hình đúng cho SoC ASR6601 và hoạt
động ổn định trong môi trường thực tế. Trước hết, nhóm sử dụng mã nguồn ví dụ
\texttt{pingpong} được cung cấp trong ASR6601 SDK làm nền tảng phát triển, sau đó
tiến hành tùy chỉnh và mở rộng chức năng để xây dựng phiên bản
\texttt{pingpong\_uart}. Phiên bản này được bổ sung cơ chế giao tiếp UART, cho phép
module LoRa trao đổi dữ liệu trực tiếp với vi điều khiển trung tâm ESP32.

Sau khi hoàn tất việc chỉnh sửa mã nguồn, project firmware được tạo tự động thông qua
công cụ \texttt{KeilProjectGen}. Script này có nhiệm vụ sinh file dự án Keil với các
thiết lập phù hợp cho kiến trúc ASR6601, giúp giảm thiểu các sai sót trong quá trình
cấu hình thủ công. Project sau đó được mở bằng môi trường Keil $\mu$Vision~5 để
thực hiện quá trình biên dịch. Khi quá trình build hoàn tất thành công, hệ thống sẽ
tạo ra file nhị phân \texttt{pingpong.bin}, là đầu vào cho bước nạp firmware.

Ở bước cuối cùng, firmware được nạp vào bộ nhớ flash của module Ra-08H thông qua
phần mềm TremoProg. Để đưa module vào chế độ nạp, chân BOOT0 được kéo lên mức
logic cao (3.3V) và module được reset. Sau đó, người vận hành lựa chọn đúng cổng
COM, chỉ định file \texttt{.bin} tương ứng và thực hiện lệnh nạp. Khi quá trình này
kết thúc thành công, module sẵn sàng hoạt động với firmware LoRa P2P đã được triển
khai.

\subsection{Quá trình phát triển firmware LoRa P2P}
Nhóm đã phát triển hai phiên bản firmware riêng biệt cho Node gửi và Node nhận dựa trên giao thức LoRa P2P (Peer-to-Peer).

\subsubsection{Node TX (Truyền tin)}
Firmware của Node TX được xây dựng với mục tiêu đóng vai trò như một bộ chuyển đổi
dữ liệu từ giao tiếp UART sang kênh truyền không dây LoRa (UART-to-LoRa). Trong
giai đoạn khởi tạo, cổng UART1 của module Ra-08H được cấu hình với tốc độ truyền
115200~bps nhằm đảm bảo khả năng giao tiếp ổn định và phù hợp với vi điều khiển
ESP32. Việc lựa chọn baudrate này giúp cân bằng giữa tốc độ truyền dữ liệu và độ
ổn định của hệ thống trong điều kiện vận hành thực tế.

Để xử lý dữ liệu nhận được từ ESP32, firmware sử dụng cơ chế lưu trữ tạm thời thông
qua bộ đệm. Kích thước bộ đệm \texttt{BUFFER\_SIZE} được mở rộng lên 256~bytes
nhằm đảm bảo có thể chứa đầy đủ chuỗi dữ liệu dạng JSON được gửi từ các cảm biến
trong mỗi chu kỳ đo. Cách tiếp cận này giúp hạn chế hiện tượng tràn bộ đệm và đảm
bảo tính toàn vẹn của dữ liệu trước khi tiến hành truyền không dây.

Về mặt logic hoạt động, vi xử lý liên tục giám sát dữ liệu trên cổng UART. Khi phát
hiện ký tự kết thúc dòng (\texttt{\textbackslash n}), firmware xác định rằng một gói
dữ liệu hoàn chỉnh đã được nhận. Toàn bộ nội dung trong bộ đệm sau đó được đóng
gói thành payload LoRa và được truyền đi bằng cách gọi hàm
\texttt{Radio.Send()}. Cơ chế này cho phép Node TX thực hiện truyền dữ liệu một cách
tuần tự, đảm bảo mỗi gói tin LoRa chứa đầy đủ thông tin cần thiết trước khi phát sóng.
Đoạn mã dưới đây trích xuất từ file \texttt{pingpong\_tx.c}, minh họa logic chính trong vòng lặp vô hạn của Node TX. Vi xử lý liên tục giám sát dữ liệu trên cổng UART, đọc từng byte vào bộ đệm \texttt{line}. Khi phát hiện ký tự kết thúc dòng (\texttt{\textbackslash n}) hoặc khi bộ đệm đầy, firmware sẽ gọi hàm \texttt{Radio.Send()} để phát sóng gói tin LoRa đi:

\begin{lstlisting}[language=C, caption={Logic đọc UART và gửi LoRa tại Node TX}, label={lst:tx_logic}]
    while (1) {
        while (uart_available_wrapper() && idx < (int)sizeof(line)-1) {
            int c = uart_readbyte_wrapper();
            if (c < 0) break;

#if UART0_ECHO
            uart_writebyte((uint8_t)c);   
#endif
            line[idx++] = (uint8_t)c;
            last_rx_tick = TimerGetCurrentTime();
            if (c == '\n') break;
        }

        if (idx > 0) {
            int have_nl = (line[idx-1] == '\n');
            int idle    = (TimerGetElapsedTime(last_rx_tick) > 80);

            if (have_nl || (idle && idx > 0)) {
                Radio.Send(line, (uint8_t)idx);
                printf("Sent %d bytes\r\n", idx);
                
                /* Reset buffer */
                idx = 0;
                memset(line, 0, sizeof(line));
            }
        }
        Radio.IrqProcess();
    }
\end{lstlisting}

\subsubsection{Node RX (Gateway nhận tin)}
Firmware của Node RX được thiết kế để hoạt động như một Gateway trung gian, thực
hiện chức năng chuyển đổi dữ liệu từ kênh truyền không dây LoRa sang giao tiếp
UART (LoRa-to-UART). Node RX đóng vai trò là điểm thu thập dữ liệu từ các Node TX
trong hệ thống và chuyển tiếp thông tin này tới vi điều khiển trung tâm hoặc máy
tính để xử lý ở các tầng cao hơn.

Khi module LoRa nhận được một gói tin hợp lệ, sự kiện nhận được kích hoạt thông qua
hàm callback \texttt{OnRxDone} do thư viện Radio cung cấp. Tại thời điểm này,
firmware tiến hành trích xuất payload dữ liệu gốc từ gói tin nhận được, đồng thời thu
thập các thông số đặc trưng của kênh truyền bao gồm cường độ tín hiệu thu (RSSI) và
tỷ lệ tín hiệu trên nhiễu (SNR). Các thông số này phản ánh chất lượng liên kết vô
tuyến và đóng vai trò quan trọng trong việc đánh giá hiệu năng truyền thông của
hệ thống.

Sau khi hoàn tất quá trình xử lý và làm giàu dữ liệu, firmware tiến hành đóng gói
payload LoRa cùng với các giá trị RSSI và SNR thành một chuỗi dữ liệu dạng JSON.
Chuỗi JSON này sau đó được gửi ra cổng UART để chuyển tiếp tới ESP32 hoặc máy
tính phục vụ cho các bước xử lý tiếp theo, lưu trữ và hiển thị dữ liệu trong hệ thống.
Cách tiếp cận này cho phép tách biệt rõ ràng giữa tầng truyền thông không dây và
tầng xử lý ứng dụng, đồng thời tạo điều kiện thuận lợi cho việc mở rộng và phân tích
hiệu năng hệ thống trong các giai đoạn tiếp theo.
Đoạn mã sau thể hiện quá trình xử lý khi nhận được gói tin (State = RX). Dữ liệu payload được mã hóa sang Hex để đảm bảo an toàn, sau đó được đóng gói cùng với RSSI, SNR, tần số và timestamp thành một chuỗi JSON hoàn chỉnh trước khi gửi qua UART:

\begin{lstlisting}[language=C, caption={Logic đóng gói JSON và gửi UART tại Node RX}, label={lst:rx_logic}]
switch (State)
{
    case RX: {
        char hex[(BUFFER_SIZE*2)+1];
        hex_encode(RxBuf, RxLen, hex);

        char json[BUFFER_SIZE*3 + 128];
        int n = snprintf(json, sizeof(json),
            "{\"ts\":%lu,\"freq\":%lu,\"rssi\":%d,\"snr\":%d,"
            "\"len\":%u,\"payload_hex\":\"%s\"}\n",
            (unsigned long)TimerGetCurrentTime(),
            (unsigned long)RF_FREQUENCY,
            (int)RxRssi, (int)RxSnr,
            (unsigned)RxLen, hex);
        if (n > 0) uart0_write((const uint8_t*)json, n);

        Radio.Rx(RX_TIMEOUT_VALUE);
        State = LOWPOWER;
        break;
    }
    
    case RX_TIMEOUT:
    case RX_ERROR:
        Radio.Rx(RX_TIMEOUT_VALUE);
        State = LOWPOWER;
        break;
    // ...
}
\end{lstlisting}
Cách tiếp cận này cho phép tách biệt rõ ràng giữa tầng truyền thông không dây và tầng xử lý ứng dụng, đồng thời tạo điều kiện thuận lợi cho việc mở rộng và phân tích hiệu năng hệ thống trong các giai đoạn tiếp theo.

% ============================================================
% 5.3 Cài đặt backend
% ============================================================
\section{Cài đặt backend và xử lý dữ liệu}

\subsection{Thiết lập hạ tầng MQTT Broker}
Hệ thống sử dụng Mosquitto làm MQTT Broker đóng vai trò trung tâm trong việc định tuyến và phân phối dữ liệu. Để đảm bảo khả năng truy cập từ các thiết bị ngoại vi trong mạng nội bộ, cấu hình của Broker đã được tùy chỉnh để lắng nghe trên tất cả các giao diện mạng (0.0.0.0) tại cổng tiêu chuẩn 1883, thay vì chỉ giới hạn ở localhost như mặc định. Đồng thời, các quy tắc tường lửa (Inbound Rules) trên máy chủ Windows cũng được thiết lập để cho phép các kết nối TCP đến cổng này, đảm bảo luồng dữ liệu từ Gateway đến Broker được thông suốt.

\subsection{Triển khai dịch vụ cầu nối (Bridge Service)}
Để kết nối giữa luồng dữ liệu thời gian thực từ MQTT và cơ sở dữ liệu lưu trữ lâu dài Firebase, nhóm đã phát triển một dịch vụ trung gian (Bridge Service) bằng ngôn ngữ Python. Dịch vụ này không chỉ đơn thuần là chuyển tiếp dữ liệu mà còn thực hiện chức năng chuẩn hóa và làm giàu dữ liệu (Data Enrichment).

Cụ thể, script Python thực hiện đăng ký (subscribe) các chủ đề (topic) dữ liệu từ Gateway thông qua wildcard \texttt{ammo\_transport/+/telemetry}. Khi nhận được gói tin, script tiến hành phân tích cú pháp JSON, trích xuất các trường thông tin quan trọng và thực hiện chuyển đổi định dạng. Dữ liệu thô từ cảm biến được kết hợp với các thông số kỹ thuật của mạng LoRa (RSSI, SNR) để tạo thành bản ghi hoàn chỉnh.

Dữ liệu sau xử lý được ghi vào Cloud Firestore theo mô hình hai lớp: lớp \textit{Latest State} phục vụ cho việc hiển thị trạng thái tức thời trên Dashboard, và lớp \textit{History} lưu trữ chuỗi thời gian (time-series) phục vụ cho nhu cầu truy xuất lịch sử và phân tích xu hướng. Việc sử dụng thư viện \texttt{firebase-admin} với cơ chế xác thực qua Service Account đảm bảo tính bảo mật và quyền kiểm soát truy cập dữ liệu nghiêm ngặt.

% ============================================================
% 5.4 Cài đặt ứng dụng di động
% ============================================================
\section{Cài đặt và phát triển ứng dụng di động}

\subsection{Kiến trúc ứng dụng Flutter}
Ứng dụng di động được xây dựng dựa trên Flutter SDK, áp dụng kiến trúc phân tầng rõ ràng để tách biệt giữa giao diện người dùng (UI), logic nghiệp vụ (Business Logic) và các dịch vụ nền tảng (Services). Mã nguồn được tổ chức thành các module chức năng riêng biệt:
\begin{itemize}
    \item \textbf{Module Giao diện (Screens \& Components):} Bao gồm các màn hình tương tác chính và các widget tái sử dụng. Các thành phần này chỉ chịu trách nhiệm hiển thị dữ liệu và tiếp nhận thao tác người dùng, không chứa logic xử lý phức tạp.
    \item \textbf{Module Dịch vụ (Services):} Tập trung các lớp xử lý kết nối mạng, quản lý trạng thái và tương tác với API bên ngoài. Điển hình là \texttt{MqttService}, chịu trách nhiệm duy trì kết nối bền vững với Broker và xử lý các sự kiện ngắt kết nối hoặc lỗi mạng.
    \item \textbf{Module Mô hình dữ liệu (Models):} Định nghĩa cấu trúc dữ liệu ánh xạ từ các bản tin JSON, giúp đảm bảo tính nhất quán và an toàn kiểu dữ liệu (type safety) trong toàn bộ ứng dụng.
\end{itemize}

\subsection{Cơ chế đồng bộ dữ liệu thời gian thực}
Thay vì sử dụng phương pháp hỏi vòng (polling) truyền thống gây tốn băng thông và năng lượng, ứng dụng áp dụng cơ chế lắng nghe sự kiện (event-driven) thông qua giao thức MQTT. Dịch vụ \texttt{MqttService} đăng ký theo dõi các chủ đề dữ liệu liên quan, cho phép ứng dụng nhận được cập nhật ngay lập tức khi có gói tin mới từ thiết bị giám sát. Dữ liệu nhận về được phân tích và cập nhật trực tiếp vào trạng thái ứng dụng (State Management), kích hoạt việc vẽ lại giao diện (re-build UI) để hiển thị thông tin mới nhất cho người dùng với độ trễ thấp nhất.

% ============================================================
% 5.5 Kết nối và tích hợp phần cứng
% ============================================================
% \section{Kết nối và tích hợp phần cứng}
% Việc tích hợp phần cứng giữa vi điều khiển ESP32 và module truyền thông Ra-08H đóng vai trò then chốt trong việc đảm bảo độ tin cậy của luồng dữ liệu. Sơ đồ kết nối vật lý được thiết kế tối ưu hóa cho giao tiếp UART không đồng bộ, cụ thể như Bảng \ref{tab:wiring}.

% \begin{table}[H]
%     \centering
%     \caption{Sơ đồ kết nối phần cứng giữa ESP32 và Module Ra-08H}
%     \label{tab:wiring}
%     \begin{tabular}{|c|c|l|}
%     \hline
%     \textbf{Chân ESP32} & \textbf{Chân Ra-08H} & \multicolumn{1}{c|}{\textbf{Chức năng và Ý nghĩa kỹ thuật}} \\ \hline
%     GPIO 17 (TX2) & RX & Kênh truyền dữ liệu chính (Data Uplink) từ vi xử lý xuống module. \\ \hline
%     GPIO 16 (RX2) & TX & Kênh nhận phản hồi (Feedback) và thông tin debug từ module LoRa. \\ \hline
%     GND & GND & Tham chiếu điện áp chung, bắt buộc để đảm bảo mức logic chính xác. \\ \hline
%     3.3V / 5V & VCC & Cung cấp năng lượng hoạt động ổn định cho module RF. \\ \hline
%     \end{tabular}
% \end{table}

% Trong quá trình triển khai thực tế, nhóm đặc biệt chú trọng đến vấn đề cấp nguồn, đảm bảo dòng điện đủ lớn để đáp ứng các đỉnh tiêu thụ (peak current) khi module LoRa thực hiện phát sóng, tránh hiện tượng sụt áp gây reset vi điều khiển.
% ============================================================
% 5.5 Kết nối và tích hợp thiết bị
% ============================================================
\section{Kết nối và tích hợp thiết bị}

Hệ thống bao gồm hai khối phần cứng chính: Node TX (Sensor Node) gắn trên phương tiện vận chuyển và Node RX (Gateway) đặt tại trạm giám sát. Cả hai đều sử dụng vi điều khiển ESP32 làm bộ xử lý trung tâm và module Ra-08H để truyền thông LoRa P2P.

\subsection{Sơ đồ kết nối Node TX (Sensor Node)}
Node TX chịu trách nhiệm thu thập dữ liệu từ đa cảm biến và gửi đi. Dựa trên mã nguồn firmware, sơ đồ đấu nối giữa ESP32 và các module ngoại vi được mô tả chi tiết trong Bảng \ref{tab:wiring_tx}.

\begin{table}[h]
    \centering
    \begin{tabular}{|p{3cm}|p{2.5cm}|p{3cm}|p{4cm}|}
    \hline
    \textbf{Module} & \textbf{Giao tiếp} & \textbf{Chân ESP32} & \textbf{Chân Module} \\ \hline
    \multirow{2}{*}{LoRa Ra-08H} & \multirow{2}{*}{UART1} & GPIO 26 (TX) & RX \\ \cline{3-4} 
     & & GPIO 25 (RX) & TX \\ \hline
    \multirow{2}{*}{GPS NEO-6M} & \multirow{2}{*}{UART2} & GPIO 17 (TX) & RX \\ \cline{3-4} 
     & & GPIO 16 (RX) & TX \\ \hline
    DHT22 & GPIO & GPIO 14 & DATA \\ \hline
    \multirow{2}{*}{ADXL345} & \multirow{2}{*}{I2C} & GPIO 21 (SDA) & SDA \\ \cline{3-4} 
     & & GPIO 22 (SCL) & SCL \\ \hline
    Nguồn & Power & 3.3V / GND & VCC / GND \\ \hline
    \end{tabular}
    \caption{Sơ đồ chân kết nối phần cứng Node TX}
    \label{tab:wiring_tx}
\end{table}



\textbf{Lưu ý cấu hình trong Firmware:}
\begin{itemize}
    \item \textbf{LoRa UART:} Sử dụng \texttt{HardwareSerial(1)} với tốc độ Baudrate là \textbf{115200} để đảm bảo đồng bộ với firmware của Ra-08H.
    \item \textbf{GPS UART:} Sử dụng tốc độ \textbf{9600} baud tiêu chuẩn của module NEO-6M.
    \item \textbf{Cảm biến gia tốc ADXL345:} Sử dụng giao tiếp I2C mặc định của ESP32 (Wire library).
\end{itemize}

\subsection{Sơ đồ kết nối Node RX (Gateway Node)}
Node RX đóng vai trò nhận dữ liệu LoRa và chuyển tiếp. Cấu hình phần cứng tối giản hơn, chỉ bao gồm ESP32 kết nối với module LoRa Ra-08H để thực hiện chức năng Gateway.

\begin{table}[h]
    \centering
    \begin{tabular}{|p{3cm}|p{2.5cm}|p{3cm}|p{4cm}|}
    \hline
    \textbf{Module} & \textbf{Giao tiếp} & \textbf{Chân ESP32} & \textbf{Chân Module} \\ \hline
    \multirow{2}{*}{LoRa Ra-08H} & \multirow{2}{*}{UART1} & GPIO 26 (TX) & RX \\  
     & & GPIO 25 (RX) & TX \\ \hline
    Máy tính (Server) & UART0 (USB) & USB Port & USB Port \\ \hline
    \end{tabular}
    \caption{Sơ đồ chân kết nối phần cứng Node RX}
    \label{tab:wiring_rx}
\end{table}


\subsection{Tích hợp dữ liệu (Data Integration)}
Để đảm bảo tính toàn vẹn dữ liệu khi truyền giữa ESP32 và Ra-08H, hệ thống sử dụng định dạng chuỗi \textbf{JSON}. Trong tác vụ \texttt{TaskLoraSend} của ESP32, các dữ liệu cảm biến rời rạc được đóng gói thành một bản tin duy nhất trước khi gửi qua UART tới module LoRa.

Cấu trúc bản tin JSON được định nghĩa như sau:
\begin{lstlisting}[language=C, caption={Đoạn code đóng gói JSON trên Node TX}]
snprintf(payload, sizeof(payload),
  "{\"device_id\":\"Ra-08H-Node1\","
  "\"timestamp\":%lu,"
  "\"temp\":%.2f,"
  "\"battery\":%.2f,"
  "\"gps\":{\"lat\":%.4f,\"lng\":%.4f},"
  "\"rssi_lora\":%d}\n",
  timestamp, temp, battery, lat, lng, rssi);
\end{lstlisting}

Việc sử dụng JSON giúp Node RX dễ dàng giải mã (parse) và Gateway có thể đẩy trực tiếp dữ liệu này lên MQTT Broker mà không cần xử lý định dạng phức tạp.
% ============================================================
% 5.6 Quy trình vận hành hệ thống
% ============================================================
\section{Quy trình vận hành và kiểm thử hệ thống}
Hệ thống được vận hành theo quy trình khởi động tuần tự từ trung tâm ra ngoại vi (Center-to-Edge) để đảm bảo tính sẵn sàng của hạ tầng trước khi thiết bị đầu cuối hoạt động.

\textbf{Bước 1: Khởi động hạ tầng máy chủ.}
Quá trình bắt đầu bằng việc kích hoạt dịch vụ Mosquitto MQTT Broker và chạy script Python Bridge. Tại bước này, các kết nối đến cơ sở dữ liệu Firebase được thiết lập và xác thực, đồng thời các cổng lắng nghe mạng được mở để sẵn sàng tiếp nhận kết nối.

\textbf{Bước 2: Kích hoạt Gateway (Node RX).}
Gateway sau khi được cấp nguồn sẽ thực hiện quy trình tự kiểm tra (POST), kết nối vào mạng WiFi cục bộ và thiết lập phiên làm việc với MQTT Broker. Sau khi kết nối thành công, Gateway chuyển sang chế độ lắng nghe liên tục trên kênh tần số LoRa đã định.

\textbf{Bước 3: Vận hành thiết bị giám sát (Node TX).}
Thiết bị giám sát gắn trên phương tiện vận tải được khởi động cuối cùng. Vi điều khiển ESP32 bắt đầu chu kỳ thu thập dữ liệu từ các cảm biến GPS, nhiệt độ và gia tốc kế. Dữ liệu được xử lý, đóng gói và truyền đi định kỳ qua module LoRa.

\textbf{Bước 4: Giám sát và Điều hành.}
Người dùng cuối truy cập vào hệ thống thông qua ứng dụng di động. Tại đây, luồng dữ liệu từ hiện trường được trực quan hóa thành các thông số giám sát và vị trí trên bản đồ số. Hệ thống cho phép người điều hành theo dõi trạng thái vận chuyển theo thời gian thực và nhận các cảnh báo tức thời khi có bất thường xảy ra.

\begin{table}[H]
    \centering
    \begin{tabular}{|l|p{10cm}|}
        \hline
        \textbf{Thành phần} & \textbf{Vai trò kỹ thuật và Chức năng} \\ \hline
        Node TX & \textbf{Thu thập biên (Edge Acquisition):} Đo đạc thông số môi trường và vận hành, xử lý tín hiệu sơ bộ và truyền tải không dây tầm xa. \\ \hline
        Node RX & \textbf{Cổng kết nối (Gateway):} Chuyển đổi môi trường truyền tin từ sóng vô tuyến (RF) sang mạng IP, đóng vai trò điểm tập trung dữ liệu. \\ \hline
        MQTT Broker & \textbf{Trục truyền tin (Message Bus):} Đảm bảo phân phối dữ liệu tin cậy, thời gian thực giữa các thành phần phân tán. \\ \hline
        Python Bridge & \textbf{Xử lý trung gian (Middleware):} Chuẩn hóa dữ liệu, tích hợp logic nghiệp vụ và đồng bộ hóa với hệ thống lưu trữ đám mây. \\ \hline
        Firestore & \textbf{Lưu trữ đám mây (Cloud Storage):} Quản lý dữ liệu bền vững, hỗ trợ truy vấn lịch sử và đồng bộ trạng thái đa thiết bị. \\ \hline
        Mobile App & \textbf{Giao diện người dùng (HMI):} Trực quan hóa dữ liệu giám sát, cung cấp công cụ tương tác và ra quyết định cho người quản lý. \\ \hline
    \end{tabular}
    \caption{Tổng hợp chức năng và vai trò của các thành phần trong hệ thống}
    \label{tab:system_components}
\end{table}
